% ══════════════════════════════════════════════════════════════
%  Style Sheet --- Visual reference for every element type
% ══════════════════════════════════════════════════════════════
\chapter{Visual Style Sheet}
\label{ch:stylesheet}

\epigraph{The root problem with conventional currency is all the trust
that's required to make it work.}{---Satoshi Nakamoto, 2009}

\begin{chapteroutline}
  \item Section Heading (Level 2)
  \item Code Listings
  \item Mathematics
  \item Tables
  \item Callout Boxes
  \item Lists
  \item Additional Elements
\end{chapteroutline}

\lettrine{T}{his} page demonstrates every visual element used in the book.
It serves as a design reference during the \LaTeX{} conversion and
ensures that all typographic choices work harmoniously together.

% ── Section headings ─────────────────────────────────────────
\section{Section Heading (Level 2)}

Body text in TeX Gyre Pagella (Palatino family) at 11pt with 1.05
line spread. Microtype provides character protrusion and font
expansion for optically even margins. Inline code appears as
\code{BlockHeader} or \code{Vec<Transaction>}---note the subtle
background pill that makes code pop from prose. Cross-references
use \cref{ch:stylesheet}.\marginnote{Cross-refs are clickable in PDF}

A \keyterm{UTXO} (Unspent Transaction Output) is the fundamental
unit of spendable value. First use of a key term prints bold and
adds an index entry automatically.

\subsection{Subsection Heading (Level 3)}

Subsections use TeX Gyre Heros (Helvetica family) in the
\textcolor{sectioncolor}{dark teal} accent color.

\subsubsection{Subsubsection Heading (Level 4)}

The deepest heading level, still sans-serif and colored.

% ── Code listing ─────────────────────────────────────────────
\section{Code Listings}

{\small\sffamily\textbf{Listing: Block structure definition (\code{block.rs})}}
\label{lst:block-struct}
\refstepcounter{listing}
\begin{minted}[fontsize=\footnotesize]{rust}
pub struct BlockHeader {     // @\cmark{1}@
    pub timestamp: i64,
    pub pre_block_hash: String,
    pub hash: String,
    pub nonce: i64,
    pub height: usize,
}

pub struct Block {           // @\cmark{2}@
    pub header: BlockHeader,
    pub transactions: Vec<Transaction>,
}
\end{minted}

\begin{codeannotations}
  \annotate{1}{The header contains only the fields needed for
    proof-of-work hashing and chain linking.}
  \annotate{2}{A block pairs its header with a vector of transactions,
    keeping structure and payload separate.}
\end{codeannotations}

A code listing without line numbers (for shell commands):

\begin{minted}[fontsize=\footnotesize]{bash}
cargo build --release
cargo test -- --nocapture
docker compose up -d
\end{minted}

% ── Inline math ──────────────────────────────────────────────
\section{Mathematics}

Inline math: the probability that an attacker catches up from
$z$ blocks behind is $q_z = (q/p)^z$ where $p$ is the probability
of an honest node finding the next block and $q = 1 - p$.

% ── Tables ───────────────────────────────────────────────────
\section{Tables}

\label{tab:wp-types}
\medskip
{\small\sffamily\textbf{Table: Whitepaper concept to Rust type mapping}}
\smallskip

\begin{center}
\small
\rowcolors{2}{white}{tableshade}
\begin{tabularx}{0.9\textwidth}{lll}
\toprule
\textbf{Whitepaper Concept} & \textbf{Rust Type} & \textbf{Chapter} \\
\midrule
Block header       & \code{BlockHeader}  & 9 \\
Transaction        & \code{Transaction}  & 9 \\
Input (spend)      & \code{TXInput}      & 9 \\
Output (new coin)  & \code{TXOutput}     & 9 \\
UTXO set           & \code{UTXOSet}      & 9 \\
\bottomrule
\end{tabularx}
\end{center}

% ── Callout boxes ────────────────────────────────────────────
\section{Callout Boxes}

All 16 canonical callout types:

\begin{notebox}
This is a \textbf{Note} callout. Used for supplementary information
that clarifies or expands on the main text.
\end{notebox}

\begin{tipbox}
This is a \textbf{Tip} callout. Practical advice the reader can
apply immediately.\marginnote{Tips are green-accented}
\end{tipbox}

\begin{warnbox}
This is a \textbf{Warning} callout. Highlights a common mistake
or non-obvious behavior.
\end{warnbox}

\begin{importantbox}
This is an \textbf{Important} callout. Critical information the
reader must not overlook.
\end{importantbox}

\begin{infobox}{Methods involved}
\code{BlockchainService::add\_block()},
\code{UTXOSet::update()},
\code{ProofOfWork::run()}
\end{infobox}

\begin{infobox}{Source}
\code{bitcoin/src/chain/blockchain\_service.rs} --- BlockchainService
\end{infobox}

\begin{learnbox}
\begin{itemize}
  \item How blocks link to form a chain via hash pointers
  \item The UTXO model for tracking spendable outputs
  \item How proof-of-work secures the chain
\end{itemize}
\end{learnbox}

\begin{prereqbox}
\begin{itemize}
  \item Chapters~\ref{ch:primitives}--\ref{ch:cryptography} (Primitives, Utilities, Cryptography)
  \item Basic familiarity with hash functions and digital signatures
\end{itemize}
\end{prereqbox}

\begin{examplebox}{See it in action}
Run \code{cargo test test\_new\_block} to see block creation and
PoW in action.
\end{examplebox}

\begin{examplebox}{Implementation example}
The \code{add\_block()} method validates the block, updates the
chain tip, and triggers a UTXO set refresh.
\end{examplebox}

\begin{refbox}{See the full implementation}
\code{bitcoin/src/chain/blockchain\_service.rs} ---
full source with inline comments.
\end{refbox}

\begin{refbox}{Companion Chapter}
Complete source code is available in Chapter~9A: Code Listings.
In the print edition, these listings appear in the Appendix:
Source Reference.
\end{refbox}

\begin{infobox}{Scope}
This section covers the domain model only. State management,
persistence, and consensus are covered in Sections~9.2--9.7.
\end{infobox}

\begin{checkpointbox}
At this point you should understand: (1) why blocks contain a
\code{pre\_block\_hash} field, (2) how transactions encode the
UTXO worldview, and (3) where the signature lives in our model.
\end{checkpointbox}

\begin{troublebox}
If \code{cargo test} fails with ``database locked,'' ensure no
other node instance is running. The Sled database uses exclusive
file locks.
\end{troublebox}

\begin{examplebox}{Framework Comparison}
Iced uses the MVU (Model--View--Update) pattern; Tauri splits
the Rust backend from a React/TypeScript frontend.
\end{examplebox}

% ── Lists ────────────────────────────────────────────────────
\section{Lists}

Bulleted list:
\begin{itemize}
  \item \code{pre\_block\_hash}: link to the previous block (whitepaper \S3/\S4)
  \item \code{nonce} + \code{hash}: PoW search result (whitepaper \S4)
  \item \code{timestamp}: ordering signal; also part of the PoW header bytes
\end{itemize}

Numbered list:
\begin{enumerate}
  \item Define what a transaction \emph{is} by reading the struct fields.
  \item Show how a spending transaction is \emph{constructed} in code.
  \item Trace the UTXO selection, signing, and broadcast flow.
\end{enumerate}

\scenbreak

% ── Additional elements ─────────────────────────────────────
\section{Additional Elements}

The scene break above (\texttt{\textbackslash scenbreak}) creates a
visual pause between major topics within a chapter.

The \texttt{\textbackslash keyterm} command marks first definitions:
a \keyterm{block header} is the metadata that identifies a block
and links it to the previous one. Terms are indexed automatically.

Epigraphs appear at the start of each chapter (see the Satoshi
quote at the top of this chapter). They set the tone without
interrupting the technical flow.

Drop caps (\texttt{\textbackslash lettrine}) open each chapter's
first paragraph with a large decorative initial in the brand color,
a hallmark of premium book typography.

Margin notes\marginnote{Like this one!}\ appear in the outer margin
for quick cross-references and asides without interrupting flow.

The \texttt{\textbackslash cmark} command creates numbered annotation
circles for code walkthrough explanations (see \cref{lst:block-struct}).

% ── Pull quote ───────────────────────────────────────────────
\pullquote{The UTXO set is derived state---not history. It answers
one question: what is spendable right now?}

% ── Deep dive sidebar ───────────────────────────────────────
\begin{deepdive}{Why we chose Sled over LevelDB}
Our implementation uses the Sled embedded database instead of
LevelDB (which Bitcoin Core uses). Sled is written in pure Rust,
provides lock-free concurrent access, and integrates naturally
with Rust's ownership model. The trade-off is a larger on-disk
footprint and less battle-tested production history. For a
learning project, the ergonomic benefits outweigh these costs.
\end{deepdive}

% ── Step-by-step tutorial ────────────────────────────────────
A tutorial walkthrough with distinct step badges:

\begin{steps}
  \step{Clone the repository}
    Run \code{git clone https://github.com/bkunyiha/\allowbreak rust-blockchain}
    to get a local copy of the project.

  \step{Build in release mode}
    Run \code{cargo build --release} from the project root.
    This compiles with optimizations enabled.

  \step{Run the test suite}
    Verify everything works: \code{cargo test -- --nocapture}.
\end{steps}

% ── Footnote demo ────────────────────────────────────────────
Footnotes\footnote{Like this one---note the rust-orange marker
and thin brand-colored separator rule above.} use brand-colored
markers for visual consistency with the rest of the design system.

% ── API reference table ──────────────────────────────────────
\medskip
{\small\sffamily\textbf{Table: API Reference --- \code{BlockchainService}}}

\smallskip
\begin{center}
\small
\rowcolors{2}{white}{tableshade}%
\begin{tabular}{@{} l l l p{4.8cm} @{}}
\toprule
\textbf{Method} & \textbf{Parameters} & \textbf{Returns} & \textbf{Description} \\
\midrule
  \code{add\_block} & \texttt{block: Block} & \texttt{Result<()>} & Validates and appends a block \\
  \code{get\_tip} & \texttt{---} & \texttt{Option<Block>} & Returns the chain tip \\
  \code{get\_block} & \texttt{hash: \&str} & \texttt{Option<Block>} & Retrieves a block by hash \\
\bottomrule
\end{tabular}
\end{center}

% ── Chapter-end summary ──────────────────────────────────────
\begin{chaptersummary}
  \item All visual elements render correctly in Crown Quarto format
  \item Callout boxes have icon prefixes for instant recognition
  \item Code listings support annotation callouts with numbered circles
  \item Pull quotes, deep dives, and step environments add variety
  \item Tables use alternating row shading; API tables have a standard format
  \item Drop caps, margin notes, footnotes, and epigraphs add polish
  \item A colophon page documents the book's production details
\end{chaptersummary}
